\subsection{\problemblock{*Define} Data Block}\label{sec:problem-block-define}

TAOS always computes a standard set of variables for each trajectory. Additional
variables can be defined within a trajectory definition that are functions of these
standard variables (\cref{sec:trajectory-block-define}). These user-defined
variables are defined only within a trajectory. If a quantity of interest is a function
of variables from more than one trajectory, it must be defined outside of the
trajectory definitions. Then, trajectory numbers can be used to reference values
from different trajectories.

The difference between \problemblock{*define} data blocks within a trajectory
definition and those outside of a trajectory definition is that
\problemblock{*define} data blocks within a trajectory are restricted to be
functions of that trajectory's output variables, and data blocks outside a trajectory
are not restricted to be functions of a single trajectory.
\problemblock{*define} data blocks outside of trajectory definitions can be used to
compute relative values comparing two or more trajectories.

User-defined variables are input the same way regardless of whether the
\problemblock{*define} data block is within a trajectory definition or not. A unique
name is given followed by equations that show how to compute the new variable.
User-defined variables that are defined outside of trajectory definitions cannot be
integral variables; this type of user-defined variable is only allowed within a
trajectory definition.

The equations defining how to compute the user-defined variable are given in a
simplified C language which is similar to Fortran. For example,

\begin{lstlisting}[style=taosmanual]
*define del_alt
  del_alt = (alt[1] - alt[2]) / 3280.84;
\end{lstlisting}

computes the difference in altitude between trajectory numbers 1 and 2 in
kilometers. The name of the new variable is \statevar{del_alt}. It is given following
the \problemblock{*define} keyword, and it must also appear on the left side of an
equation somewhere in the definition. In this example, with one equation,
\statevar{del_alt} must be the variable on the left side of the equal sign.

The values and variables on the right side of the equal sign must all be standard
output variables, previously input user-defined variables, temporary variables, or
constants. In this case, \statevar{alt} is a standard output variable and 3280.84 is a
constant. Since this data block is located outside of the trajectory definitions in the
problem file, trajectory numbers are required to fully specify the variables.
Trajectory numbers are enclosed in square brackets so they look like subscripts.

\Cref{tab:user-defined-math} contains a list of operators and
functions that can be used in equations for user-defined variables. All of the standard
mathematical operators are available, and they have the same precedence as in C and
Fortran. Many standard functions are also provided; again these are defined in TAOS
as they are in the C language.

The \texttt{table()} function shown in \cref{tab:user-defined-math}
cannot be used in \problemblock{*define} data blocks outside of a trajectory
definition. Tables are functions of the state of a vehicle or trajectory so they must be
associated with a specific trajectory. They cannot be evaluated outside of a
trajectory.

Equations must end with a semicolon like in C. The semicolon is required even when
there is only one equation like in the example shown above. Since equations continue
until a semicolon is encountered, they can extend across multiple lines as necessary.

When more than one equation is required to define a value, temporary variables can
be used. These variables must be unique, that is, they cannot have the same name as
a standard output variable or another user-defined variable. They are only defined
within the \problemblock{*define} data block. For example,

\begin{lstlisting}[style=taosmanual]
*define rel_range
  dx = (east[2] - east[1])^2;
  dy = (north[2] - north[1])^2;
  rel_range = sqrt(dx+dy);
\end{lstlisting}

computes a relative horizontal range between two trajectories using two temporary
variables, \statevar{dx} and \statevar{dy}, in the calculation.

If-then-else statements can be used as shown in
\cref{sec:trajectory-block-define}, but \texttt{while} loops, \texttt{for} loops,
and other C-language control statements are not allowed.

User-defined variables given outside of the trajectory definitions can only be
referenced by \problemblock{*print}, \problemblock{*file},
\problemblock{*egs}, and other \problemblock{*define} data blocks that are also
outside of the trajectory definitions.
